% Tapis de Quillen, with Quillen's offprint and preprint (cote 162-5, typed part): archive cote 162-5, site page 46 (PDF page 47).
% typescript. Draft compared: none.
% AI-transcribed by Claude Opus 5.5 (claude-opus-5-5, adaptive thinking, effort high), Message Batches API, batch t162_full (msgbatch_01UAh3zyWPd454AeDkhnmQ33), 2026-11-13; not yet checked by eye.
\scanpage{162-5/46}
catégorie \inserted[typed]{additive(et même quasi-abélienne)}avec $\otimes$ qui a toutes les vertus,sauf d'être abélienne, et
sauf d'admettre des \uline{Hom} internes. Pour trouver ces derniers, on peut
si on veut paraphraser la construction de la catégorie des motifs, en
\xout{\unsure{introduisant un motif}} considérant l'objet \inserted{$\Lambda(-1)$} de $\tilde{\underline{B}}$ conoyau de
$p^{\hand{*}}$ : $e \hand{$\longrightarrow$} S$, où $p:S^{1} \hand{$\longrightarrow$} e$ est l'application du cercle sur un
point. C'est un objet de poids 1 (le poids est défini par le degré \xout{d}
des générateurs, changés de signe). On constate aussitôt que le
foncteur $M\hand{$\mapsto$} M(-1)=M\otimes\hand{$\Lambda$}(-1)$ de $\tilde{\underline{B}}$ dans lui-même est pleinement fidèle,
ce qui permet alors, par une construction formelle standard, de complé-%
ter $\overset{\hand{$\sim$}}{\underline{B}}$ en une catégorie où $\hand{$\Lambda$}(-1)$ devient inversible. Il se trouve
alors que dans cette dernière les \uline{Hom}(X,Y) internes existent et
\xout{satisfont} sont isomorphes à $\check{X}\otimes Y$ ... en fait , on trouve maintenant la
catégorie de\xout{s} \uline{tous} les modules gradués libres de type fini sur \quad .
Bien sûr, cette catégorie n'est pas encore abélienne pour autant,
puisque \quad n'est pas un corps; \xout{pour} \inserted[typed]{si on veut} définir la catégorie abélienne
engendrée, on s'attend à trouver la\xout{s} catégorie \hand{/}desmodules de \deleted{type} \inserted{présentation} fin-%
sur $\hand{$\Lambda$}$.( \inserted{\unsure{qui est abélienne bien que} $\Lambda$ \unsure{ne soit pas noethérien, car} $\Lambda$ \unsure{est}} \deleted{\quad est-il noethérien, i.e. le nombre de générateurs est-il}
\margin{\unsure{?}} \deleted{fini ? ).} \hand{\unsure{limite d'anneaux noeth. à homs de transition} \uline{plats})}.

Cette construction ressemblerait à celle d'une théorie des
motifs où on n'aurait pas passé à l'équivalence homologique pour
les cycles, mais où on aurait gardé l'équivalence rationelle, - qui
est bien trop fine \xout{\unsure{pou}}à beaucoup d'égards. Pour trouver un équiva\xout{l}%
lent plus étroit, il convient de reprendre la théorie prédédente,on
pourrait formuler le même Pb universel que plus haut, mais en prenant
l'équivalence homologique (via le graphe) des morphismes de variétés,
au lieu de l'équivalence homotopique. \deleted{Mais alors il semble que l'on
trouve comme catégorie solution celle dont les objets sont les varié-}

\mspage{46}
